\subsection{逻辑函数}
逻辑函数可以表示为
\begin{equation}
    Y=F(A_1,A_2,\cdots,A_n).
\end{equation}
其中 $Y,A_1,A_2,\cdots,A_n$ 都只可取 0 或 1.

\subsubsection{逻辑函数的表示方法}
\begin{itemize}
    \item 逻辑真值表;
    \item 逻辑表达式 (具体的函数解析式);
    \item 逻辑图 (用逻辑元件的图形表示);
    \item 波形图 (时序图);
    \item 卡诺图;
    \item ...
\end{itemize}

\subsubsection{各种表示方法的转换}
以\textbf{真值表转换为逻辑表达式}为例:
\begin{enumerate}
    \item 找出真值表中使逻辑函数输出为 1 的组合;
    \item 对于每一种组合, 令为 1 的输入变量记为原变量, 令为 0 的输入变量记为其反变量, 并将所有输入变量组成乘积项;
    \item 将乘积项相加即可 (\textbf{与或式}).
\end{enumerate}

\subsubsection{逻辑函数的两种标准形式}
\begin{itemize}
    \item 标准与或式 (最小项之和);
    \item 标准或与式 (最大项之积).
\end{itemize}

\textbf{最小项}\quad $n$ 个逻辑变量的最小项满足
\begin{itemize}
    \item 是一个连续乘积项;
    \item 包含所有这 $n$ 个变量;
    \item 每一个变量都以原变量或其反变量的形式\textit{出现且仅有一次}.
\end{itemize}
显然 $n$ 个变量的最小项共有 $2^n$ 个. 当 $n=3$ 时, 这些最小项可列表如下:
\begin{table}[H]
    \caption{三变量的所有最小项}
    \centering
    \begin{tabular}{ccccc} \toprule
        $A$ & $B$ & $C$ & \bfseries 最小项 & $m_i$ \\\midrule
        0   & 0   & 0   & $A'B'C'$      & $m_0$ \\
        0   & 0   & 1   & $A'B'C$       & $m_1$ \\
        0   & 1   & 0   & $A'BC'$       & $m_2$ \\
        0   & 1   & 1   & $A'BC$        & $m_3$ \\
        1   & 0   & 0   & $AB'C'$       & $m_4$ \\
        1   & 0   & 1   & $AB'C$        & $m_5$ \\
        1   & 1   & 0   & $ABC'$        & $m_6$ \\
        1   & 1   & 1   & $ABC$         & $m_7$ \\
        \bottomrule
    \end{tabular}
\end{table}

\underline{性质}
\begin{itemize}
    \item 对于任意一个最小项, 仅有一组输入变量取值使其为 1, 其他均为 0 (在输入变量的任何取值下, 所有最小项中有且只有一个为 1, 其他均为 0);
    \item $n$ 变量组成\textit{所有}最小项之和为 1;
    \item $n$ 变量组成\textit{任二}最小项之积为 0;
    \item 具有\textbf{相邻性} (有且仅有一个因子不同) 的两个最小项之和可以合并, 并消除不同的那个因子.
\end{itemize}

\textbf{最大项}\quad 与最小项类似, 用 $M_i$ 表示.

相同标号的最大项和最小项满足如下关系
\begin{equation}
    m_i'=M_i.
\end{equation}

\textbf{标准与或式的函数式的特点}
\begin{itemize}
    \item 函数式为与或形式;
    \item 不一定包含所有最小项, 但所有项均为最小项.
\end{itemize}

标准与或式的函数式可以表示为
\begin{equation}
    Y=m_{k_1}+m_{k_2}+\cdots+m_{k_p}=\sum m(k_1,k_2,\cdots,k_p).
\end{equation}

\textbf{两种标准形式的相互转化}
\begin{equation}
    Y=\sum_{i\in\{k_1,k_2,\cdots,k_p\}}m_i=\prod_{i\notin\{k_1,k_2,\cdots,k_p\}}M_i.
\end{equation}

例如, 当 $n=3$ 时, 有 $Y=m(1,3,5)=M(0,2,4,6,7)$.
